\subsection{Real measures. Positive measures}

Let X be a locally compact space. The real vector space \(\mathcal{K}(\mathrm{X};\mathbf{R})\) is a subspace of the real vector space underlying the complex vector space \(\mathcal{K}(\mathrm{X};\mathbf{C})\) ; moreover, the mapping \((f_{1},f_{2})\mapsto f_{1}+if_{2}\) is an isomorphism of the product topological vector space \(\mathcal{K}(\mathrm{X};\mathbf{R})\times\mathcal{K}(\mathrm{X};\mathbf{R})\) onto the real topological vector space \(\mathcal{K}(\mathrm{X};\mathbf{C})\) (No.~1, Prop. 1).

For every (complex) measure \(\mu\in\mathcal{M}(\mathrm{X};\mathbf{C})\) , the restriction \(\mu_{0}\) of \(\mu\) to \(\mathcal{K}(\mathrm{X};\mathbf{R})\) is a continuous R-linear mapping of \(\mathcal{K}(\mathrm{X};\mathbf{R})\) into C; moreover, this restriction determines \(\mu\) , since if \(f=f_{1}+if_{2}\) with \(f_{1},f_{2}\) in \(\mathcal{K}(\mathrm{X};\mathbf{R})\) , then \(\mu(f)=\mu_{0}(f_{1})+i\mu_{0}(f_{2})\) . Conversely, let \(\mu_{0}\) be a continuous R-linear mapping of \(\mathcal{K}(\mathrm{X};\mathbf{R})\) into C; it is clear that the mapping

\[
f _ {1} + i f _ {2} \mapsto \mu_ {0} (f _ {1}) + i \mu_ {0} (f _ {2})
\]

is a (complex) measure on X. Thus, every measure on X may be identified with its restriction to \(\mathcal{K}(\mathrm{X};\mathbf{R})\) .

Let \(\mu\) be a measure on X. One calls conjugate measure of \(\mu\) the measure \(\overline{\mu}\) defined by \(\overline{\mu}(f) = \overline{\mu(\overline{f})}\) for every function \(f \in \mathcal{K}(\mathrm{X};\mathbf{C})\) ; for, it is clear that \(\overline{\mu}\) is a C-linear form and that it is continuous on \(\mathcal{K}(\mathrm{X};\mathbf{C})\) ; obviously \(\overline{\overline{\mu}} = \mu\) , and, for two measures \(\mu, \nu\) and two scalars \(\alpha, \beta\) in \(\mathbf{C}\) ,

\[
\overline {{(\alpha \mu + \beta \nu)}} = \overline {{\alpha}} \cdot \overline {{\mu}} + \overline {{\beta}} \cdot \overline {{\nu}}.
\]

More generally, for every function \(g \in \mathcal{C}(\mathrm{X};\mathbf{C})\) and every measure \(\mu\) on \(\mathbf{X}\) ,

\[
\overline {{{g \cdot \mu}}} = \overline {{{g}}} \cdot \overline {{{\mu}}},\tag{7}
\]

as is immediate from the definition (No.~4).

A measure \(\mu\) on X is said to be real if \(\overline{\mu} = \mu\) ; by the foregoing, it is the same to say that for every function \(f \in \mathcal{K}(X; \mathbf{R})\) , \(\mu(f)\) is a real number. If one identifies a real measure with its restriction to \(\mathcal{K}(X; \mathbf{R})\) , one can thus say that the set of real measures on X is the dual of the real locally convex space \(\mathcal{K}(X; \mathbf{R})\) ; it is a real vector space that is denoted \(\mathcal{M}(X; \mathbf{R})\) (or sometimes \(\mathcal{M}(X)\) if this does not cause any confusion). The Lebesgue measure on R is a real measure, as is the Dirac measure \(\varepsilon_{a}\) for every point \(a \in X\) . If \(g \in \mathcal{C}(X; \mathbf{R})\) and if \(\mu\) is a real measure, then so is \(g \cdot \mu\) by virtue of (7).

Let \(\mu\) be a (complex) measure on X. It follows from the preceding definition that the measures \(\mu_{1} = (\mu +\overline{\mu}) / 2\) and \(\mu_{2} = (\mu -\overline{\mu}) / 2i\) are real; they are called, respectively, the real part and the imaginary part of \(\mu\) , and they are denoted by \(\mathcal{R}\mu\) and \(\mathcal{I}\mu\) , respectively; these measures are also characterized by the fact that, for every function \(f\in \mathcal{K}(\mathbf{X};\mathbf{R})\) ,

\[
\mu_ {1} (f) = \mathcal {R} (\mu (f)), \quad \mu_ {2} (f) = \mathcal {I} (\mu (f)).
\]

Obviously

\[
\mu = \mu_ {1} + i \mu_ {2}, \quad \overline {{{{\mu}}}} = \mu_ {1} - i \mu_ {2}.
\]

The space \(\mathcal{K}(\mathrm{X};\mathbf{R})\) of continuous real-valued functions on \(\mathrm{X}\) with compact support is clearly a Riesz space for the order relation \(f\leqslant g\) . We shall say that a real measure \(\mu\) on \(\mathrm{X}\) is positive if \(\mu (f)\geqslant 0\) for every function \(f\geqslant 0\) belonging to \(\mathcal{K}(\mathrm{X};\mathbf{R})\) ; it is thus a positive linear form on the Riesz space \(\mathcal{K}(\mathrm{X};\mathbf{R})\) (Ch. II, §2, No.~1, Def. 1). Conversely:

THEOREM 1. --- Every positive linear form on the Riesz space \(\mathcal{K}(\mathrm{X};\mathbf{R})\) is a (positive) real measure on \(\mathrm{X}\) .

For, let \(\mu\) be a positive linear form on \(\mathcal{K}(\mathrm{X};\mathbf{R})\) and let \(\mathrm{K}\) be a compact subset of \(\mathrm{X}\) . There exists a continuous mapping \(f_0\) of \(\mathrm{X}\) into {[}0, 1{]}, with compact support, such that \(f_0(x) = 1\) on \(\mathrm{K}\) (No.~2, Lemma 1). For every function \(g \in \mathcal{K}(\mathrm{X},\mathrm{K};\mathbf{R})\) , we thus have \(-\| g \| f_0 \leqslant g \leqslant \| g \| f_0\) , consequently \(|\mu(g)| \leqslant \| g \| \cdot \mu(f_0)\) , which proves the theorem.

We denote by \(\mathcal{M}_{+}(X)\) the pointed convex cone of positive measures on X (or, what amounts to the same thing, the cone of positive linear forms on the Riesz space \(\mathcal{K}(X;\mathbf{R})\) ).

THEOREM 2. --- Every real measure on a locally compact space X is the difference of two positive measures.

In view of Theorem 1 and Ch. II, §2, No.~2, Th. 1, it all comes down to proving that a real measure \(\mu\) on X is a relatively bounded linear form on the Riesz space \(\mathcal{K}(\mathrm{X};\mathbf{R})\) . Let \(f\) be a continuous function \(\geqslant 0\) on X, with compact support K; the relation \(0 \leqslant g \leqslant f\) in \(\mathcal{K}(\mathrm{X};\mathbf{R})\) implies that \(\|g\| \leqslant \|f\|\) and that the support of \(g\) is contained in K. By hypothesis, there exists a number \(\mathrm{M}_{\mathrm{K}} \geqslant 0\) such that \(|\mu(h)| \leqslant \mathrm{M}_{\mathrm{K}} \cdot \|h\|\) for every function \(h \in \mathcal{K}(\mathrm{X},\mathrm{K};\mathbf{R})\) ; therefore \(|\mu(g)| \leqslant \mathrm{M}_{\mathrm{K}} \cdot \|g\| \leqslant \mathrm{M}_{\mathrm{K}} \cdot \|f\|\) , which proves the theorem.

The space \(\mathcal{M}(\mathrm{X};\mathbf{R})\) of real measures on X is thus identical with the space of relatively bounded linear forms on the Riesz space \(\mathcal{K}(\mathrm{X};\mathbf{R})\) ; we recall that in \(\mathcal{M}(\mathrm{X};\mathbf{R})\) , the order relation \(\mu\leqslant\nu\) signifies that \(\nu-\mu\) is a positive measure, or also that \(\mu(f)\leqslant\nu(f)\) for every function \(f\in\mathcal{K}_{+}(\mathrm{X})\) .

THEOREM 3. --- The space \(\mathcal{M}(\mathrm{X};\mathbf{R})\) of real measures on a locally compact space X is fully lattice-ordered.

This follows from Ch. II, §2, No.~2, Th. 1.

In conformity with the notations of Ch. II, §1, we define, for every real measure \(\mu\) on \(X\) ,

\[
\mu^ {+} = \sup (\mu , 0), \quad \mu^ {-} = \sup (- \mu , 0), \quad | \mu | = \sup (\mu , - \mu);
\]

then \(\mu = \mu^{+} - \mu^{-}\) , \(|\mu| = \mu^{+} + \mu^{-}\) and \(\inf (\mu^{+},\mu^{-}) = 0\) . Moreover, for every function \(f\in \mathcal{K}_{+}(\mathrm{X})\) ,

\[
\int f d \mu^ {+} = \sup _ {0 \leqslant g \leqslant f, g \in \mathcal {K} (\mathrm{X})} \int g d \mu\tag{8}
\]

and

\[
\int f d | \mu | = \sup _ {| g | \leqslant f, g \in \mathcal {K} (\mathrm{X})} \int g d \mu ,\tag{9}
\]

whence, in particular,

\[
\left| \int f d \mu \right| \leqslant \int | f | d | \mu |\tag{10}
\]

for every function \(f \in \mathcal{K}(\mathbf{X};\mathbf{R})\) .

This inequality is also true if \(f \in \mathcal{K}(\mathrm{X};\mathbf{C})\) ; for, on multiplying f by a complex number of absolute value 1 (which does not change either side of the inequality), one can suppose that \(\int f d\mu \geqslant 0\) . Then

\[
\left| \int f d \mu \right| = \int f d \mu = \int (\mathcal {R} f) d \mu \leqslant \int | \mathcal {R} f | d | \mu | \leqslant \int | f | d | \mu |.
\]

